\documentclass[11pt]{article}
\usepackage[margin=1in]{geometry}
\usepackage{amsmath,amssymb,amsthm}
\usepackage{hyperref}

\newtheorem{proposition}{Proposition}
\newtheorem{lemma}{Lemma}
\newtheorem{corollary}{Corollary}

\newcommand{\R}{\mathbb{R}}

\title{Exact line search step for a quadratic}
\author{Companion note \texttt{1800-03} to \texttt{1800-HagerZhang\_2013}}
\date{}

\begin{document}
\maketitle

\section{What the paper uses}

For the quadratic $\tfrac12x^{\top}Ax-b^{\top}x$, the Fletcher--Reeves code of
the paper takes the step $\alpha_k=\|g_k\|^2/d_k^{\top}Ad_k$, and Table~1.1 is
computed ``with an exact line search''~\cite[pp.~2150--2152]{HZ13}. This note
proves that these are the same thing, and gives the step used for PALMER1C.

\section{Setting}

Let
\[
f(x)=\tfrac12x^{\top}Hx-c^{\top}x+\gamma,\qquad x\in\R^n,
\]
with $H\in\R^{n\times n}$ symmetric positive definite, $c\in\R^n$, $\gamma\in\R$.
Its gradient is $g(x)=Hx-c$. For a point $x$ and a direction $d\ne0$, the
\emph{exact line search} step is the minimizer of $\varphi(\alpha):=f(x+\alpha d)$
over $\alpha\in\R$.

\section{Result}

\begin{lemma}[$f$ along a line]\label{lem}
For all $x,d\in\R^n$ and $\alpha\in\R$,
\[
\varphi(\alpha)=f(x)+\alpha\,g(x)^{\top}d+\tfrac12\alpha^2\,d^{\top}Hd .
\]
\end{lemma}

\begin{proof}
Expand the quadratic term:
\[
\tfrac12(x+\alpha d)^{\top}H(x+\alpha d)
=\tfrac12x^{\top}Hx+\tfrac12\alpha\,(x^{\top}Hd+d^{\top}Hx)+\tfrac12\alpha^2d^{\top}Hd
=\tfrac12x^{\top}Hx+\alpha\,x^{\top}Hd+\tfrac12\alpha^2d^{\top}Hd,
\]
because $d^{\top}Hx=(d^{\top}Hx)^{\top}=x^{\top}H^{\top}d=x^{\top}Hd$ ($H$
symmetric). The linear term gives $-c^{\top}(x+\alpha d)=-c^{\top}x-\alpha\,c^{\top}d$.
Adding these and $\gamma$,
\[
\varphi(\alpha)=f(x)+\alpha\,(x^{\top}H-c^{\top})d+\tfrac12\alpha^2d^{\top}Hd .
\]
Finally $(x^{\top}H-c^{\top})d=(Hx-c)^{\top}d=g(x)^{\top}d$, again by symmetry of $H$.
\end{proof}

\begin{proposition}[Exact step]\label{prop}
Let $d\ne0$, and write $p:=g(x)^{\top}d$, $q:=d^{\top}Hd$. Then $q>0$, and
$\varphi$ has the unique minimizer
\[
\alpha^\ast=-\frac{g(x)^{\top}d}{d^{\top}Hd}.
\]
Moreover $g(x+\alpha^\ast d)^{\top}d=0$, and $\alpha^\ast>0$ whenever $d$ is a
descent direction ($p<0$).
\end{proposition}

\begin{proof}
\emph{$q>0$:} $H$ is positive definite and $d\ne0$.

\emph{Minimizer.} By Lemma~\ref{lem}, $\varphi(\alpha)=f(x)+p\alpha+\tfrac12q\alpha^2$.
Completing the square with $\alpha^\ast=-p/q$:
\[
\tfrac12q(\alpha-\alpha^\ast)^2-\frac{p^2}{2q}
=\tfrac12q\alpha^2-q\alpha^\ast\alpha+\tfrac12q(\alpha^\ast)^2-\frac{p^2}{2q}
=\tfrac12q\alpha^2+p\alpha+\frac{p^2}{2q}-\frac{p^2}{2q}
=p\alpha+\tfrac12q\alpha^2 .
\]
Hence
\[
\varphi(\alpha)=\varphi(\alpha^\ast)+\tfrac12q(\alpha-\alpha^\ast)^2,\qquad
\varphi(\alpha^\ast)=f(x)-\frac{p^2}{2q}.
\]
Since $q>0$, $\varphi(\alpha)\ge\varphi(\alpha^\ast)$, with equality if and only
if $\alpha=\alpha^\ast$.

\emph{Orthogonality.} $g(x+\alpha^\ast d)=Hx-c+\alpha^\ast Hd=g(x)+\alpha^\ast Hd$, so
$g(x+\alpha^\ast d)^{\top}d=p+\alpha^\ast q=p-p=0$.

\emph{Sign.} If $p<0$ then $\alpha^\ast=-p/q>0$, since $q>0$.
\end{proof}

\section{The two cases used}

\begin{corollary}[PALMER1C]\label{cor:palmer}
Let $f(a)=\|Va-y\|^2$ with $V$ of full column rank, $A=V^{\top}V$, $b=V^{\top}y$.
Then $g(a)=2(Aa-b)$ and the exact step is
\[
\alpha^\ast=-\frac{g(a)^{\top}d}{2\,d^{\top}Ad}.
\]
\end{corollary}

\begin{proof}
$f(a)=(Va-y)^{\top}(Va-y)=a^{\top}V^{\top}Va-2y^{\top}Va+y^{\top}y
=\tfrac12a^{\top}(2A)a-(2b)^{\top}a+y^{\top}y$, using $a^{\top}V^{\top}y=y^{\top}Va$.
This is the setting with $H=2A$, $c=2b$, $\gamma=y^{\top}y$, so $g(a)=2(Aa-b)$.
$H$ is positive definite: for $v\ne0$, $v^{\top}(2A)v=2\|Vv\|^2>0$ because
$V$ has full column rank. Proposition~\ref{prop} with $d^{\top}Hd=2\,d^{\top}Ad$
gives the formula.
\end{proof}

For PALMER1C, $V_{ik}=x_i^{2k}$ ($k=0,\dots,7$) is a Vandermonde matrix in
$t_i=x_i^2$; it has full column rank because the 35 data points give 18 distinct
values of $t_i$ (at least 8 are needed).

\begin{corollary}[The paper's step]\label{cor:paper}
Let $f(x)=\tfrac12x^{\top}Ax-b^{\top}x$ with $A$ symmetric positive definite,
and run the Fletcher--Reeves code~\cite[p.~2152]{HZ13} with exact steps:
$d_0=-g_0$, $d_{k+1}=-g_{k+1}+\beta_kd_k$, and $\alpha_k$ the minimizer of
$f(x_k+\alpha d_k)$. Then for every $k$ with $g_k\ne0$,
\[
g_k^{\top}d_k=-\|g_k\|^2\qquad\text{and}\qquad
\alpha_k=\frac{\|g_k\|^2}{d_k^{\top}Ad_k},
\]
which is the step in the code.
\end{corollary}

\begin{proof}
Here $H=A$. \emph{First identity, by induction on $k$.} For $k=0$,
$g_0^{\top}d_0=-\|g_0\|^2$. For $k\ge0$, $\alpha_k$ is the exact step along
$d_k$, so $g_{k+1}^{\top}d_k=0$ by Proposition~\ref{prop}. Then
$g_{k+1}^{\top}d_{k+1}=-\|g_{k+1}\|^2+\beta_k\,g_{k+1}^{\top}d_k=-\|g_{k+1}\|^2$.

\emph{Step.} $g_k\ne0$ gives $d_k\ne0$ (as $g_k^{\top}d_k=-\|g_k\|^2\ne0$), so
Proposition~\ref{prop} applies:
$\alpha_k=-g_k^{\top}d_k/(d_k^{\top}Ad_k)=\|g_k\|^2/(d_k^{\top}Ad_k)$.
\end{proof}

The identity $g_k^{\top}d_k=-\|g_k\|^2$ relies on all previous steps being
exact. If one step is not exact (as in Example~2 of note \texttt{1800-02}),
the formula $\|g_k\|^2/d_k^{\top}Ad_k$ is no longer the exact step, while
$-g_k^{\top}d_k/d_k^{\top}Ad_k$ still is.

\begin{thebibliography}{9}
\bibitem{HZ13} W.~W.~Hager and H.~Zhang, The limited memory conjugate gradient
method, \emph{SIAM J. Optim.}, 23 (2013), pp.~2150--2168.
\end{thebibliography}

\end{document}
